\documentclass[11pt,a4paper]{article}
\usepackage[T1]{fontenc}
\usepackage[utf8]{inputenc}
\usepackage{lmodern}
\usepackage[french]{babel}
\usepackage[margin=2cm,headheight=14pt]{geometry}
\usepackage[most]{tcolorbox}
\usepackage{xcolor,tabularx,booktabs,array,enumitem,fancyhdr,amssymb,tikz}
\usetikzlibrary{arrows.meta,positioning}

\pagestyle{fancy}
\fancyhf{}
\lhead{\small Robotique -- Grade 6 -- Novembre 2026}
\rhead{\small Leçon 4 : Choisir un matériau adapté}
\cfoot{\thepage}
\setlist{itemsep=1pt,topsep=2pt,leftmargin=*}
\renewcommand{\arraystretch}{1.35}
\newcolumntype{Y}{>{\raggedright\arraybackslash}X}
\newcommand{\cb}{$\square$\,}
\newcommand{\fl}{$\rightarrow$}
\newcommand{\lig}{\noindent\rule{\linewidth}{0.3pt}\\[5mm]}

\newtcolorbox{etape}[3]{enhanced,breakable,colback=#1!6,colframe=#1!70!black,
  coltitle=white,fonttitle=\bfseries,title={#2\hfill #3},boxrule=0.8pt,arc=2mm,
  before skip=8pt,after skip=8pt}
\newtcolorbox{astuce}[1][Différenciation et besoins particuliers]{breakable,
  colback=yellow!12,colframe=yellow!50!black,fonttitle=\bfseries\small,
  title={#1},boxrule=0.5pt,arc=1mm,before skip=6pt,after skip=6pt}
\newcommand{\entete}[2]{%
  \begin{tcolorbox}[colback=#1!10,colframe=#1!70!black,boxrule=0.8pt,arc=2mm]
  {\Large\bfseries #2}\hfill Nom : \rule{4cm}{0.3pt}\quad Date : \rule{2.2cm}{0.3pt}
  \end{tcolorbox}}

% Schéma de la démarche de choix (Besoin -> Critères -> Comparaison -> Choix)
\newcommand{\demarche}{%
\begin{center}
\begin{tikzpicture}[box/.style={draw,rounded corners,thick,minimum height=1.3cm,text width=3cm,align=center,font=\small\bfseries},>={Stealth[length=3mm]},node distance=0.8cm]
\node[box,fill=orange!20] (a) {1. Besoin\\[-1mm]\footnotesize\mdseries à quoi sert l'objet ?};
\node[box,fill=blue!15,right=of a] (b) {2. Critères\\[-1mm]\footnotesize\mdseries que faut-il ?};
\node[box,fill=green!15,right=of b] (c) {3. Comparer\\[-1mm]\footnotesize\mdseries tableau ++ / + / $-$};
\node[box,fill=violet!15,right=of c] (d) {4. Choisir\\[-1mm]\footnotesize\mdseries et justifier};
\draw[->,thick] (a)--(b); \draw[->,thick] (b)--(c); \draw[->,thick] (c)--(d);
\end{tikzpicture}
\end{center}}

\begin{document}

\begin{tcolorbox}[colback=teal!10,colframe=teal!70!black,arc=2mm,boxrule=1pt]
\centering
{\LARGE\bfseries Leçon 4 : Choisir un matériau adapté à un objet technique}\\[2mm]
{\large Robotique -- Grade 6 -- Plan de cours de 2 heures}\\[1mm]
Enseignant : Francis Ki \quad|\quad Semaine du 9 au 13 novembre 2026
\end{tcolorbox}

%==================================================================
\section{Informations générales}
\begin{tabularx}{\linewidth}{>{\bfseries}p{3.6cm}Y}
\toprule
Titre & Leçon 4 : Choisir un matériau adapté à un objet technique\\
Classe / effectif & Grade 6, 11 à 20 élèves, niveaux hétérogènes, élèves à besoins particuliers\\
Durée & 2 heures (120 min)\\
Lieu & Salle informatique, un poste par élève\\
Matériel et logiciels &
\begin{itemize}
 \item Un poste par élève avec navigateur web à jour et accès internet.
 \item Vidéoprojecteur ou tableau interactif pour l'enseignant.
 \item Diapositives (Google Slides ou LibreOffice Impress) pour la pratique du niveau C ; traitement de texte ou tableur pour les tableaux comparatifs des niveaux A et B.
 \item Fiches numériques partagées en lecture seule : chaque élève crée sa copie et la renomme avec son prénom.
 \item Facultatif : trois cartons de vote par élève (vert, rouge) pour l'accroche, ou un sondage en ligne.
\end{itemize}\\
Pré-requis &
Leçon 3 : connaître les quatre familles de matériaux (bois, métal, plastique, verre) et au moins deux propriétés de chacune. Savoir utiliser la phrase modèle \textit{«~La\ldots{} est en\ldots{} car elle est\ldots~»}.\\
\bottomrule
\end{tabularx}

\begin{astuce}[Préparation avant le cours (environ 30 min)]
\begin{enumerate}
 \item Préparer les fichiers numériques à partir des annexes B à E (une copie par élève).
 \item Préparer cinq phrases projetées pour l'accroche (voir étape 1).
 \item Créer le formulaire de sortie (annexe G) dans Google Forms ou imprimer l'annexe.
 \item Reconstituer les trois groupes de besoin à partir de l'exit ticket de la leçon 3.
\end{enumerate}
\end{astuce}

%==================================================================
\section{Objectifs d'apprentissage}
\begin{tcolorbox}[colback=gray!8,colframe=gray!60!black,arc=1mm,boxrule=0.5pt]
\textbf{Objectifs du programme}
\begin{itemize}
 \item Comparer différents matériaux selon leurs caractéristiques.
 \item Justifier le choix d'un matériau en fonction d'un besoin.
\end{itemize}
\textbf{À la fin de la séance, l'élève est capable\,:}
\begin{itemize}
 \item d'énoncer le besoin auquel répond un objet et d'en déduire deux ou trois critères\,;
 \item de comparer deux ou trois matériaux dans un tableau avec les symboles ++, + et $-$\,;
 \item de rédiger un choix justifié : \textit{«~Je choisis\ldots{} parce que\ldots, même si\ldots~»}.
\end{itemize}
\textbf{Compétences numériques\,:} remplir un tableau comparatif, mettre en forme une diapositive, enregistrer et rendre son travail.
\end{tcolorbox}

%==================================================================
\section{Déroulé minuté (120 min)}

\begin{tabularx}{\linewidth}{>{\bfseries}p{2.2cm}p{2.6cm}Y p{2.8cm}}
\toprule
Minutes & Étape & Contenu & Organisation\\
\midrule
0 -- 10 & Accroche & Vrai ou absurde ? & Classe entière, sans écran\\
10 -- 30 & Mini-leçon & Les quatre temps du choix d'un matériau & Classe entière\\
30 -- 60 & Pratique guidée & Une gourde collective, puis deux objets en binômes & Collectif puis binômes\\
60 -- 95 & Pratique autonome & Trois niveaux de besoin & Individuel\\
95 -- 110 & Mise en commun & Valoriser chaque groupe & Classe entière\\
110 -- 120 & Exit ticket & Trois questions sur poste & Individuel\\
\bottomrule
\end{tabularx}

%------------------------------------------------------------------
\begin{etape}{orange}{1. Accroche : vrai ou absurde ?}{0 -- 10 min}
\textbf{But\,:} faire sentir qu'un matériau ne se choisit pas au hasard.\\
\textbf{Écrans éteints.} Projeter une phrase à la fois. Les élèves votent par carton vert (c'est logique) ou rouge (c'est absurde), puis un élève justifie.
\begin{enumerate}
 \item Une casserole en bois.
 \item Une fenêtre en métal.
 \item Un casque de vélo en verre.
 \item Un pneu en bois.
 \item Une bouteille d'eau en plastique.
\end{enumerate}
Question-problème : \textbf{«~Comment fait-on pour choisir le bon matériau ?~»} Noter trois idées au tableau, qui serviront à construire les critères.
\end{etape}

%------------------------------------------------------------------
\begin{etape}{blue}{2. Mini-leçon : la démarche de choix}{10 -- 30 min}
\textbf{Support\,:} schéma construit au tableau (annexe A). Quatre temps de 5 min.

\demarche

\textbf{Temps 1 : le besoin.} À quoi sert l'objet, où et par qui est-il utilisé ? Exemple : une gourde pour le sport.

\textbf{Temps 2 : les critères.} Un critère est une qualité demandée au matériau. On en retient deux ou trois, par exemple léger, ne casse pas, sans danger pour la santé.

\textbf{Temps 3 : comparer.} Le tableau suivant, qui sert de référence, est projeté puis commenté. Les notes sont des ordres de grandeur.

\begin{center}
\begin{tabularx}{\linewidth}{>{\bfseries}p{4.2cm}Y Y Y Y}
\toprule
Critère & Bois & Métal & Plastique & Verre\\
\midrule
Léger & ++ & $-$ & ++ & $-$\\
Résiste aux chocs & + & ++ & + & $-$\\
Isole (électricité, chaleur) & ++ & $-$ & ++ & +\\
Transparent & $-$ & $-$ & + & ++\\
Résiste à l'eau & $-$ & + & ++ & ++\\
Résiste à la chaleur & $-$ & ++ & $-$ & +\\
Prix peu élevé & + & $-$ & ++ & +\\
Facile à recycler & + & ++ & $-$ & ++\\
\bottomrule
\end{tabularx}
\end{center}
\textit{Légende : ++ très bien, + correct, $-$ peu adapté.}

\textbf{Temps 4 : choisir et justifier.} Compter les ++ qui répondent \emph{aux critères choisis} (et non à tous). Rédiger : \textit{«~Je choisis\ldots{} parce que\ldots, même si\ldots~»}. Un choix est un compromis, il n'y a pas toujours un matériau parfait.

\begin{astuce}
Garder le schéma des quatre étapes affiché toute la séance. Associer chaque critère à un pictogramme. Pour les élèves lents en lecture, lire le tableau à voix haute et utiliser le lecteur immersif.
\end{astuce}
\end{etape}

%------------------------------------------------------------------
\begin{etape}{green}{3. Pratique guidée}{30 -- 60 min}
\textbf{Fiche\,:} annexe B (version numérique : fichier «~L4\_Guidee~»).

\medskip
\textbf{3a. Exemple collectif sur poste : une gourde de sport (12 min).} Chaque élève suit l'enseignant sur son poste.
\begin{enumerate}
 \item Besoin : boire pendant le sport, dans le sac.
 \item Choisir trois critères parmi la liste du tableau de référence.
 \item Comparer plastique, métal et verre en recopiant les notes.
 \item Choisir et rédiger la phrase de justification.
\end{enumerate}

\textbf{3b. Deux objets en binômes (18 min).} Cas 1 : le manche d'une casserole. Cas 2 : la coque d'un casque de vélo. Alternance des rôles à mi-temps.

\textbf{Étayage progressif de l'aide.}
\begin{center}
\begin{tabularx}{\linewidth}{>{\bfseries}p{2.4cm}Y p{4.6cm}}
\toprule
Niveau d'aide & Contenu & Qui intervient\\
\midrule
Aide 1 & Tableau de référence (annexe A) & Le binôme seul\\
Aide 2 & Questions guides : \emph{Quel est le besoin ? Quel matériau serait dangereux ou impossible ? Pourquoi ?} & Le binôme avec la carte\\
Aide 3 & L'enseignant remplit une ligne avec le binôme, qui termine le tableau & L'enseignant\\
\bottomrule
\end{tabularx}
\end{center}

\begin{astuce}
Écrire la phrase à trous au tableau pour les élèves qui ont du mal à rédiger. Accepter un schéma annoté ou une réponse orale en complément de l'écrit.
\end{astuce}
\end{etape}

%------------------------------------------------------------------
\begin{etape}{violet}{4. Pratique autonome différenciée}{60 -- 95 min}
\textbf{Consigne commune\,:} chaque élève ouvre le fichier de son groupe, le renomme avec son prénom et travaille seul. L'enseignant passe d'abord auprès du groupe A pendant 10 min.

\medskip
\begin{tabularx}{\linewidth}{>{\bfseries}p{3cm}Y p{3.2cm}}
\toprule
Groupe & Tâche sur poste & Annexe\\
\midrule
A : en dessous du niveau &
Trois objets, avec un matériau et une raison à relier dans une banque de mots. Phrases de justification à trous. Guidage visuel fort. & Annexe C\\
B : au niveau &
Trois objets (cuillère à café, fenêtre de serre, banc de jardin) : établir seul le tableau comparatif, choisir et justifier avec «~parce que\ldots{} même si\ldots~». & Annexe D\\
C : au-dessus du niveau &
Défi de conception d'une boîte à repas : cahier des charges avec des contraintes qui se contredisent, choix d'un matériau par pièce, diapositive de présentation. & Annexe E\\
\bottomrule
\end{tabularx}

\medskip
\textbf{Plan B si un élève termine tôt\,:} proposer la fiche du niveau supérieur ou le défi \textit{«~Quel objet de la classe aurait pu être fait d'un autre matériau, et qu'est-ce que cela changerait ?~»}
\end{etape}

%------------------------------------------------------------------
\begin{etape}{teal}{5. Mise en commun}{95 -- 110 min}
\begin{enumerate}
 \item (2 min) Retour au calme, écrans en veille.
 \item (9 min) Trois passages courts de 3 min, avec valorisation :
 \begin{itemize}
  \item \textbf{Groupe A} : lecture d'un choix justifié. Féliciter l'emploi de «~parce que~».
  \item \textbf{Groupe B} : présentation d'un tableau comparatif. Valoriser la rigueur de la comparaison.
  \item \textbf{Groupe C} : présentation de la diapositive de la boîte à repas. Valoriser la gestion du compromis.
 \end{itemize}
 \item (4 min) Revenir à la question-problème et compléter ensemble : \textit{On choisit un matériau en partant du besoin, puis des critères.}
\end{enumerate}
\end{etape}

\begin{etape}{gray}{6. Clôture : exit ticket}{110 -- 120 min}
Formulaire en ligne ou papier (annexe G) :
\begin{enumerate}
 \item Pourquoi choisit-on du métal pour une casserole ?
 \item Complète : Je choisis \ldots{} parce que \ldots, même si \ldots
 \item Auto-évaluation : \emph{Je maîtrise / Je progresse / J'ai besoin d'aide.}
\end{enumerate}
\end{etape}

%==================================================================
% ANNEXES
%==================================================================
\newpage
\section*{Annexe A : Support résumé de la mini-leçon}
\entete{blue}{Choisir un matériau adapté}

\demarche

\begin{tcolorbox}[colback=yellow!10,colframe=yellow!50!black]
\textbf{Un critère} est une qualité que le matériau doit avoir pour que l'objet fonctionne.\\
\textbf{Un choix} est souvent un compromis : on garde le matériau qui répond le mieux aux critères les plus importants.
\end{tcolorbox}

\begin{tabularx}{\linewidth}{>{\bfseries}p{4.2cm}Y Y Y Y}
\toprule
Critère & Bois & Métal & Plastique & Verre\\
\midrule
Léger & ++ & $-$ & ++ & $-$\\
Résiste aux chocs & + & ++ & + & $-$\\
Isole (électricité, chaleur) & ++ & $-$ & ++ & +\\
Transparent & $-$ & $-$ & + & ++\\
Résiste à l'eau & $-$ & + & ++ & ++\\
Résiste à la chaleur & $-$ & ++ & $-$ & +\\
Prix peu élevé & + & $-$ & ++ & +\\
Facile à recycler & + & ++ & $-$ & ++\\
\bottomrule
\end{tabularx}

\medskip
\begin{tcolorbox}[colback=green!8,colframe=green!50!black]
\textbf{Phrase modèle :} Je choisis \underline{\hspace{2.5cm}} parce que \underline{\hspace{3cm}}, même si \underline{\hspace{3cm}}.
\end{tcolorbox}

%==================================================================
\newpage
\section*{Annexe B : Fiche de pratique guidée}
\entete{green}{Pratique guidée : choisir un matériau}

\subsection*{3a. Exemple collectif : une gourde de sport}
Besoin : \dotfill\\
Mes trois critères : \cb léger \quad \cb ne casse pas \quad \cb isolant \quad \cb transparent \quad \cb recyclable

\begin{tabularx}{\linewidth}{>{\bfseries}p{3.8cm}Y Y Y}
\toprule
Critère & Plastique & Métal & Verre\\
\midrule
 & & & \\ & & & \\ & & & \\
\bottomrule
\end{tabularx}

\medskip
\textbf{Je choisis} \dotfill{} \textbf{parce que} \dotfill, \textbf{même si} \dotfill.

\subsection*{3b. En binôme : deux objets}
\textbf{Cas 1 : le manche d'une casserole.} Besoin : \dotfill

\begin{tabularx}{\linewidth}{>{\bfseries}p{3.8cm}Y Y Y}
\toprule
Critère & Bois & Métal & Plastique\\
\midrule
 & & & \\ & & & \\
\bottomrule
\end{tabularx}
Choix et justification : \dotfill

\medskip
\textbf{Cas 2 : la coque d'un casque de vélo.} Besoin : \dotfill

\begin{tabularx}{\linewidth}{>{\bfseries}p{3.8cm}Y Y Y}
\toprule
Critère & Plastique & Métal & Verre\\
\midrule
 & & & \\ & & & \\
\bottomrule
\end{tabularx}
Choix et justification : \dotfill

%==================================================================
\newpage
\section*{Annexe C : Niveau A (en dessous du niveau)}
\entete{orange}{Niveau A : je choisis avec de l'aide}

\textbf{Banque de mots :} plastique, métal, verre, bois ; léger, transparent, résistant, isole la chaleur, ne casse pas.

\subsection*{Activité 1 : relie chaque objet à son matériau et à sa raison (glisse les étiquettes)}
\begin{tabularx}{\linewidth}{>{\bfseries}p{4cm}Y Y}
\toprule
Objet & Matériau choisi & Raison\\
\midrule
Une fenêtre & & \\
Une poêle (le fond) & & \\
Une bouteille d'eau à emporter & & \\
\bottomrule
\end{tabularx}

\subsection*{Activité 2 : complète les phrases}
\begin{enumerate}
 \item Pour la fenêtre, je choisis le \dotfill{} parce qu'il est \dotfill.
 \item Pour la poêle, je choisis le \dotfill{} parce qu'il résiste à \dotfill.
 \item Pour la bouteille, je choisis le \dotfill{} parce qu'il est \dotfill{} et qu'il ne \dotfill{} pas.
\end{enumerate}

\subsection*{Activité 3 : vrai ou faux ? (coche)}
\begin{enumerate}
 \item Une casserole en plastique fondrait sur le feu. \quad \cb Vrai \quad \cb Faux
 \item Un casque de vélo en verre protégerait bien la tête. \quad \cb Vrai \quad \cb Faux
\end{enumerate}

%==================================================================
\newpage
\section*{Annexe D : Niveau B (au niveau)}
\entete{blue}{Niveau B : je compare et je justifie}

Pour chaque objet : écris le besoin, choisis trois critères, complète un tableau avec ++, + et $-$ (appuie-toi sur l'annexe A), puis écris ton choix avec «~parce que\ldots{} même si\ldots~».

\subsection*{Objet 1 : une cuillère à café pour boisson chaude}
Besoin : \dotfill

\begin{tabularx}{\linewidth}{>{\bfseries}p{3.8cm}Y Y Y Y}
\toprule
Critère & Bois & Métal & Plastique & Verre\\
\midrule
 & & & & \\ & & & & \\ & & & & \\
\bottomrule
\end{tabularx}
Choix : \dotfill

\subsection*{Objet 2 : le toit d'une serre pour plantes}
Besoin : \dotfill

\begin{tabularx}{\linewidth}{>{\bfseries}p{3.8cm}Y Y Y Y}
\toprule
Critère & Bois & Métal & Plastique & Verre\\
\midrule
 & & & & \\ & & & & \\ & & & & \\
\bottomrule
\end{tabularx}
Choix : \dotfill

\subsection*{Objet 3 : un banc de jardin}
Besoin : \dotfill

\begin{tabularx}{\linewidth}{>{\bfseries}p{3.8cm}Y Y Y Y}
\toprule
Critère & Bois & Métal & Plastique & Verre\\
\midrule
 & & & & \\ & & & & \\ & & & & \\
\bottomrule
\end{tabularx}
Choix : \dotfill

%==================================================================
\newpage
\section*{Annexe E : Niveau C (au-dessus du niveau)}
\entete{violet}{Niveau C : défi de conception, la boîte à repas}

\begin{tcolorbox}[colback=violet!6,colframe=violet!60!black]
\textbf{Cahier des charges de l'école.} Une boîte à repas doit être : \textbf{légère}, \textbf{très résistante aux chocs}, \textbf{peu chère} et \textbf{écologique}. Attention : certaines contraintes se contredisent. À toi de faire un compromis.
\end{tcolorbox}

\subsection*{Étape 1 : choisis un matériau pour chaque pièce}
\begin{tabularx}{\linewidth}{>{\bfseries}p{3cm}Y Y}
\toprule
Pièce & Matériau choisi & Justification (critères et compromis)\\
\midrule
Corps de la boîte & & \\
Couvercle & & \\
Fermeture (clips) & & \\
\bottomrule
\end{tabularx}

\subsection*{Étape 2 : quelle contrainte a été sacrifiée, et pourquoi ?}
\lig \lig

\subsection*{Étape 3 : une diapositive pour la classe}
Crée une diapositive avec : un schéma de la boîte légendé, ton tableau de choix en trois lignes, une phrase de conclusion \textit{«~Je choisis\ldots{} parce que\ldots, même si\ldots~»}.

\textbf{Défi bonus :} change une contrainte (par exemple «~très chaud~» pour un repas chaud) et dis quelles pièces tu changerais.

%==================================================================
\newpage
\section*{Annexe F : Corrections et attendus (enseignant)}

\subsection*{Accroche}
1. Absurde (le bois brûle). 2. Absurde (le métal conduit la chaleur et n'est pas transparent). 3. Absurde (le verre se casse). 4. Absurde. 5. Logique.

\subsection*{Annexe B : pratique guidée}
\begin{tabularx}{\linewidth}{>{\bfseries}p{3cm}Y Y}
\toprule
Cas & Choix attendu & Justification\\
\midrule
Gourde & Plastique ou métal (inox) & Léger, ne casse pas ; le verre est écarté car fragile\\
Manche de casserole & Bois ou plastique résistant à la chaleur & Il ne doit pas conduire la chaleur ; le métal est écarté sauf s'il est isolé\\
Coque de casque & Plastique & Léger, absorbe les chocs ; le verre casse, le métal est trop lourd\\
\bottomrule
\end{tabularx}

\subsection*{Annexe C : niveau A}
Fenêtre : verre, transparent. Poêle : métal, résiste à la chaleur. Bouteille : plastique, léger et ne casse pas. Activité 3 : Vrai, Faux.

\subsection*{Annexe D : niveau B}
Cuillère à café : métal (résiste à la chaleur) ou bois ou plastique dur selon l'usage. Toit de serre : verre (transparent, résiste à l'eau), ou plastique plus léger et moins cher. Banc : bois ou métal (résistant) ; plastique possible, le verre est à écarter.\\
\textbf{Critères de réussite :} besoin énoncé, trois critères pertinents, tableau cohérent avec l'annexe A, choix justifié avec «~même si~».

\subsection*{Annexe E : niveau C}
Exemple : corps en plastique (léger, peu cher), couvercle en plastique, fermeture en plastique ou métal si recyclage important. Sacrifié : écologique (le plastique est difficile à recycler) ou résistance.\\
\textbf{Critères de réussite :} un choix par pièce, une contrainte sacrifiée expliquée, une diapositive lisible.

%==================================================================
\newpage
\section*{Annexe G : Exit ticket et grille d'observation}
\entete{gray}{Exit ticket : choisir un matériau}

\begin{enumerate}
 \item Pourquoi choisit-on du métal pour une casserole ?\\ \lig
 \item Complète : Je choisis \dotfill{} parce que \dotfill, même si \dotfill.
 \item Je me sens : \cb Je maîtrise \quad \cb Je progresse \quad \cb J'ai besoin d'aide
\end{enumerate}

\textbf{Attendus :} 1. Il résiste à la chaleur et la conduit. 2. Phrase complète avec un choix, une raison et un inconvénient.

\bigskip
\subsection*{Grille d'observation (enseignant)}
\begin{tabularx}{\linewidth}{Y p{2.4cm} p{2.4cm} p{2.4cm} Y}
\toprule
Élève & Besoin et critères & Comparaison & Justification & Groupe suivant\\
\midrule
 & & & & \\ & & & & \\ & & & & \\ & & & & \\ & & & & \\
\bottomrule
\end{tabularx}

\end{document}
